\section{Introduction}
\label{sec:introduction}

Fine-tuning adapts pretrained language models to downstream tasks and domains through task-specific supervision, while parameter-efficient methods such as Low-Rank Adaptation (LoRA) make this adaptation possible with only a small number of trainable parameters \cite{hu2022lora}. Instruction tuning further shows that supervised adaptation across collections of tasks can improve generalization to unseen tasks \cite{wei2022finetuned,chung2024scaling}. However, an important question remains unclear: when a model learns a structured decision policy through fine-tuning, do the resulting changes represent transferable decision operations, constraint-sensitive behavioral adaptation, or merely surface-level improvements specific to the training domain?

This question is particularly important for decision tasks in which the preferred action depends on changing conditions. In realistic scenarios, effective decision making requires more than mapping an input state to a fixed output. A previously preferred action may become invalid after an additional rule, prohibition, or contextual constraint is introduced, requiring the model to suppress that action and select an alternative. We refer to this setting as \emph{rule-constrained fine-tuning}. Compared with ordinary task fine-tuning, rule-constrained examples explicitly supervise how decisions should change when the admissible action space is modified, providing a controlled environment for studying conditional decision behavior.

A key challenge in studying such transfer is separating task-specific surface forms from underlying decision operations. Improvements on conventional benchmarks may reflect better instruction following, memorized task knowledge, or shared linguistic patterns, rather than transfer of reusable reasoning structures. Therefore, a stricter evaluation should examine whether behaviors learned in one structured domain remain effective when the original terminology and context are replaced by different but structurally related tasks.

To investigate this problem, we use Dou Dizhu as a structured source domain. The choice of Dou Dizhu is not intended to study game playing itself, but to provide a compact environment that combines several properties relevant to controlled decision analysis. Unlike many classical rule-based environments that focus primarily on fixed planning or adversarial search, Dou Dizhu contains hierarchical relations, combinatorial patterns, legal action constraints, and conditional action selection within a unified decision process. These characteristics make it suitable for studying whether fine-tuning can transfer abstract decision operations beyond source-domain surface forms.

We construct a mixed fine-tuning corpus containing 50,400 general instruction examples and 5,500 Dou Dizhu decision examples. The general component is included to reduce over-specialization to the game-specific format, while the Dou Dizhu component supplies repeated structured decision signals. We evaluate two adaptation regimes built around this mixture. We refer to the ordinary condition as Dou Dizhu LoRA (DD-LoRA) and the rule-constrained condition as Rule-Constrained LoRA (RC-LoRA). DD-LoRA uses ordinary Dou Dizhu decision examples, whereas RC-LoRA modifies 851 of the 5,500 Dou Dizhu examples (15.5\%) by adding explicit prohibitions that invalidate an otherwise preferred action and require selection of a legal alternative. This three-way comparison among Base, DD-LoRA, and RC-LoRA separates the transfer associated with ordinary structured-task adaptation from the additional performance changes observed under constraint-enriched supervision.

We study this problem from three complementary perspectives. First, we evaluate whether decision structures learned from the source domain transfer to tasks with different surface forms. To this end, we introduce the Abstracted Cognitive Logic Benchmark (ACL Bench), a 300-item benchmark that abstracts source-domain decision operations into three reasoning categories: Hierarchy Logic, Pattern Recognition, and Counter Logic. ACL Bench removes domain-specific terminology while preserving underlying structural operations, enabling evaluation of whether fine-tuning is associated with transfer at the level of reasoning structure.

Second, we examine how fine-tuning changes model behavior when explicit constraints invalidate previously preferred actions. Rather than measuring only whether a prohibited action becomes less likely, we quantify how an explicit prohibition changes the relative log-probability margin between that action and the strongest legal alternative. We therefore construct controlled constraint-sensitive decision probes and introduce the Best Constraint Delta Margin (BCDM), a log-probability-based metric for measuring constraint-induced relative preference shifts.

Third, behavioral evaluation cannot reveal how the internal components supporting the measured preference shift differ between model states. For Qwen2.5-7B, we therefore perform causal intervention analysis to compare functional dependence across attention heads between Base and RC-LoRA. This analysis characterizes the internal support for the BCDM-based behavior rather than directly explaining the source of cross-task transfer.

\begin{figure}[ht]
\centering
\begin{tikzpicture}[
node distance=1.2cm,
box/.style={
draw,
rounded corners,
align=center,
minimum width=3.2cm,
minimum height=0.8cm
},
smallbox/.style={
draw,
rounded corners,
align=center,
minimum width=2.8cm,
minimum height=0.7cm
},
arrow/.style={
->,
thick
}
]


% Training data

\node (general) [smallbox]
{General Instruction\\Data};

\node (doudizhu) [smallbox, right=1.2cm of general]
{Dou Dizhu Decision\\Data};


\node (corpus) [box, below=1cm of $(general)!0.5!(doudizhu)$]
{Mixed Fine-tuning\\Corpus};


\draw[arrow] (general) -- (corpus);
\draw[arrow] (doudizhu) -- (corpus);


% Fine tuning

\node (finetune) [box, below=1cm of corpus]
{Rule-Constrained\\Fine-tuning};


\draw[arrow] (corpus) -- (finetune);


% Model

\node (model) [box, below=1cm of finetune]
{Fine-tuned\\LLM};


\draw[arrow] (finetune) -- (model);



% Three branches

\node (transfer) [box, below left=1.5cm and 1.8cm of model]
{
Task Transfer\\Analysis\\
\vspace{0.1cm}
ACL Bench\\
Hierarchy Logic\\
Pattern Recognition\\
Counter Logic
};


\node (constraint) [box, below=1.8cm of model]
{
Constraint-Sensitive\\
Preference Analysis\\
\vspace{0.1cm}
Constraint Probes\\
Best Constraint\\
Delta Margin
};


\node (mechanism) [box, below right=1.5cm and 1.8cm of model]
{
Mechanistic\\
Analysis\\
\vspace{0.1cm}
Attention-head\\
Causal Ablation\\
Top-$k$ Control
};


\draw[arrow] (model) -- (transfer);
\draw[arrow] (model) -- (constraint);
\draw[arrow] (model) -- (mechanism);


\end{tikzpicture}
\caption{
Overview of the proposed rule-constrained fine-tuning framework.
The framework integrates structured domain adaptation, cross-task transfer evaluation, constraint-sensitive behavioral analysis, and causal investigation of internal model changes.
}
\label{fig:framework}
\end{figure}


Our experiments reveal three main findings. First, ordinary DD-LoRA already produces selective cross-task transfer, but its effect is highly model-dependent: it improves overall ACL Bench performance for DeepSeek-R1-Distill-Llama-8B and Qwen2.5-7B while slightly reducing performance for Llama3.1-8B. RC-LoRA then yields a positive overall increment relative to DD-LoRA for all three checkpoints, giving a more consistently positive cross-checkpoint transfer profile under the constraint-enriched regime. Second, for Qwen2.5-7B, RC-LoRA exhibits a substantially larger BCDM than Base, indicating a larger constraint-induced relative preference shift; this probability-level change does not translate into improved rule-compliant generation and, under greedy decoding, RC-LoRA selects the prohibited action in all 60 probes. Third, the Base-to-RC-LoRA causal comparison reveals sparse and heterogeneous changes in attention-head contributions, indicating selective changes in functional dependence rather than uniform modification across the attention system.

The main contributions of this work are summarized as follows:

\begin{itemize}

    \item We establish a structured comparison among Base, DD-LoRA, and RC-LoRA, providing an empirical reference for distinguishing ordinary task-transfer effects from the additional changes associated with explicit constraint supervision.

    \item We introduce ACL Bench and BCDM to separately evaluate structural task transfer and constraint-induced relative preference shifts.

    \item We use causal interventions on attention heads to characterize changes in functional dependence associated with constraint-sensitive decision behavior after rule-constrained fine-tuning.

\end{itemize}


The remainder of this paper is organized as follows. Section~\ref{sec:related_work} reviews related work on fine-tuning transfer, constraint following, and mechanistic interpretability. Section~\ref{sec:method} introduces the proposed framework, benchmark construction, and evaluation metrics. Section~\ref{sec:experimental_setup} describes the experimental settings. Section~\ref{sec:results} presents the behavioral and causal analysis results. Section~\ref{sec:discussion} discusses limitations and implications. Finally, Section~\ref{sec:conclusion} concludes the study.
